\par\Needspace{4\baselineskip}
\begin{center}
\SourceLeaf{031b}\pdfbookmark[2]{14.}{article-014}14.
\end{center}
The equations in Article 8 provide several methods for deriving the true anomaly and the radius vector from the eccentric anomaly. We shall explain the most useful ones.

\par\Needspace{14\baselineskip}
\SourceLeaf{032}I. In the usual method, $v$ is determined by equation VII, and then $r$ by equation II. Applied to the example in the preceding article, using the value of $p$ given in Article 10, this method gives the following calculation.
\begin{center}
\begin{tabular}{@{}ll@{}}
\multicolumn{2}{l}{$\frac{1}{2}E=162^0\,8^{\prime}\,14^{\prime\prime}{,}75$}\\[.3em]
\makebox[16em][l]{$\log\operatorname{tang}\frac{1}{2}E$\hfill .\,.\,.\,.\,.\,.} & $9{,}508\,2198_n$\\[.3em]
\makebox[16em][l]{$\log\operatorname{tang}(45^0-\frac{1}{2}\varphi)$\hfill .\,.} & $9{,}891\,2427$\\[.3em]
\hline\noalign{\vskip2pt}
\makebox[16em][l]{$\log\operatorname{tang}\frac{1}{2}v$\hfill .\,.\,.\,.\,.\,.\,.} & $9{,}616\,9771_n$\\[.3em]
\multicolumn{2}{l}{$\frac{1}{2}v=157^0\,30^{\prime}\,41^{\prime\prime}{,}50$}\\[.3em]
\multicolumn{2}{l}{$v=315\ \ 1\ \ 23{,}00$}\\[.3em]
\end{tabular}
\end{center}
\begin{center}
\begin{tabular}{@{}ll@{}}
\makebox[16em][l]{$\log e$\hfill .\,.\,.\,.\,.\,.\,.\,.\,.\,.\,.\,.} & $9{,}389\,7262$\\[.3em]
\makebox[16em][l]{$\log\cos v$\hfill .\,.\,.\,.\,.\,.\,.\,.\,.\,.} & $9{,}849\,6597$\\[.3em]
\hline
 & $9{,}239\,3859$\\[.3em]
$e\cos v=$ & $0{,}173\,5345$\\[.3em]
\hline
\makebox[16em][l]{$\log p$\hfill .\,.\,.\,.\,.\,.\,.\,.\,.\,.\,.\,.} & $0{,}395\,4837$\\[.3em]
\makebox[16em][l]{$\log(1+e\cos v)$\hfill .\,.\,.\,.\,.} & $0{,}069\,4959$\\[.3em]
\hline
\makebox[16em][l]{$\log r$\hfill .\,.\,.\,.\,.\,.\,.\,.\,.\,.\,.\,.} & $0{,}325\,9878$\\[.3em]
\end{tabular}
\end{center}

II. The following method is shorter when several positions are to be calculated, since the constant logarithms of the quantities $\BareRoot{(a(1+e))}$ and $\BareRoot{(a(1-e))}$ need be computed only once. From equations V and VI we obtain
\[
\begin{gathered}
\sin\frac{1}{2}v\cdot\BareRoot{r}=\sin\frac{1}{2}E\BareRoot{(a(1+e))}\\[.5em]
\cos\frac{1}{2}v\cdot\BareRoot{r}=\cos\frac{1}{2}E\BareRoot{(a(1-e))},
\end{gathered}
\]
from which $\frac{1}{2}v$ and $\log\BareRoot{r}$ are readily determined. In general, whenever we have $P\sin Q=A$ and $P\cos Q=B$, we find $Q$ by the formula $\operatorname{tang}Q=\frac{A}{B}$, and then $P$ by $P=\frac{A}{\sin Q}$ or by $P=\frac{B}{\cos Q}$. The former is preferable when $\sin Q$ is greater than $\cos Q$, the latter when $\cos Q$ is greater than $\sin Q$. Problems that lead to equations of this kind (which will occur very often in this work) usually involve the condition that $P$ must be positive. This immediately settles whether $Q$ should be taken between $0$ and $180^0$ or between $180^0$ and $360^0$. If there is no such condition, the choice is left to us.

\par\Needspace{11\baselineskip}
In our example we have $e=0{,}245\,3162$
\begin{center}
\begin{tabular}{@{}ll@{}}
\makebox[16em][l]{$\log\sin\frac{1}{2}E$\hfill .\,.\,.\,.\,.\,.} & $9{,}486\,7632$\\[.3em]
\makebox[16em][l]{$\log\BareRoot{(a(1+e))}$\hfill .\,.\,.} & $0{,}258\,8593$\\[.3em]
\end{tabular}
\end{center}
\begin{center}
\begin{tabular}{@{}ll@{}}
\makebox[16em][l]{$\log\cos\frac{1}{2}E$\hfill .\,.\,.\,.\,.\,.} & $9{,}978\,5434_n$\\[.3em]
\makebox[16em][l]{$\log\BareRoot{(a(1-e))}$\hfill .\,.\,.} & $0{,}150\,1020$\\[.3em]
\end{tabular}
\end{center}

\SourceLeaf{033}\noindent Hence
\begin{center}
\footnotesize
\begin{tabular}{@{}ll@{}}
$\left.\begin{array}{ll}
\log\sin\frac{1}{2}v\cdot\BareRoot{r}\quad .\,. &9{,}745\,6225\\[.3em]
\log\cos\frac{1}{2}v\cdot\BareRoot{r}\quad .\,. &0{,}128\,6454_n
\end{array}\right\}$ & \begin{tabular}{@{}l@{}}
whence $\log\operatorname{tang}\frac{1}{2}v$ .\,.\,.\,. $9{,}616\,9771_n$\\[.3em]
$\frac{1}{2}v=157^0\,30^{\prime}\,41^{\prime\prime}{,}50$\\[.3em]
$v=315\ \ 1\ \ 23{,}00$
\end{tabular}\\
\begin{tabular}{@{}ll@{}}
\makebox[13em][l]{$\log\cos\frac{1}{2}v$\hfill .\,.\,.\,.\,.} & $9{,}965\,6515_n$\\[.3em]
\hline
\makebox[13em][l]{$\log\BareRoot{r}$\hfill .\,.\,.\,.\,.\,.\,.} & $0{,}162\,9939$\\[.3em]
\makebox[13em][l]{$\log r$\hfill .\,.\,.\,.\,.\,.\,.\,.} & $0{,}325\,9878$\\[.3em]
\end{tabular} &
\end{tabular}
\end{center}

III. We add a third method, which is almost as straightforward as the second. Its precision generally makes it preferable when the utmost accuracy is required. First $r$ is determined by equation III, and then $v$ by X. Our example worked out by this method is as follows.
\begin{center}
\begin{tabular}{@{}ll@{}}
\makebox[16em][l]{$\log e$\hfill .\,.\,.\,.\,.\,.\,.\,.\,.\,.\,.\,.} & $9{,}389\,7262$\\[.3em]
\makebox[16em][l]{$\log\cos E$\hfill .\,.\,.\,.\,.\,.\,.\,.\,.} & $9{,}909\,4637$\\[.3em]
\hline
 & $9{,}299\,1899$\\[.3em]
$e\cos E=$ & $0{,}199\,1544$\\[.3em]
\hline
\makebox[16em][l]{$\log a$\hfill .\,.\,.\,.\,.\,.\,.\,.\,.\,.\,.\,.} & $0{,}422\,4389$\\[.3em]
\makebox[16em][l]{$\log(1-e\cos E)$\hfill .\,.\,.\,.} & $9{,}903\,5488$\\[.3em]
\hline
\makebox[16em][l]{$\log r$\hfill .\,.\,.\,.\,.\,.\,.\,.\,.\,.\,.\,.} & $0{,}325\,9877$\\[.3em]
\end{tabular}
\end{center}
\begin{center}
\begin{tabular}{@{}ll@{}}
\makebox[16em][l]{$\log\sin E$\hfill .\,.\,.\,.\,.\,.\,.\,.\,.} & $9{,}766\,3366_n$\\[.3em]
\makebox[16em][l]{$\log\BareRoot{(1-e\cos E)}$\hfill .\,.\,.} & $9{,}951\,7744$\\[.3em]
\hline
 & $9{,}814\,5622_n$\\[.3em]
\makebox[16em][l]{$\log\sin\frac{1}{2}\varphi$\hfill .\,.\,.\,.\,.\,.\,.\,.} & $9{,}092\,0395$\\[.3em]
\hline\noalign{\vskip2pt}
\makebox[16em][l]{$\log\sin\frac{1}{2}(v-E)$\hfill .\,.\,.\,.} & $8{,}906\,6017_n$\\[.3em]
\multicolumn{2}{l}{$\frac{1}{2}(v-E)=-4^0\,37^{\prime}\,33^{\prime\prime}{,}24$}\\[.3em]
\multicolumn{2}{l}{$v-E=-9\ \ 15\ \ 6{,}48$}\\[.3em]
\multicolumn{2}{l}{$v=315\ \ 1\ \ 23{,}02$}\\[.3em]
\end{tabular}
\end{center}

\par\Needspace{10\baselineskip}
Formula VIII or XI is very convenient for checking the calculation, especially if $v$ and $r$ have been determined by the third method. The calculation is as follows.
\begin{center}
\begin{tabular}{@{}ll@{}}
\makebox[16em][l]{$\log\frac{a}{r}\sin E$\hfill .\,.\,.\,.\,.} & $9{,}862\,7878_n$\\[.3em]
\makebox[16em][l]{$\log\cos\varphi$\hfill .\,.\,.\,.\,.\,.\,.} & $9{,}986\,5224$\\[.3em]
\hline
 & $9{,}849\,3102_n$\\[.3em]
\makebox[16em][l]{$\log\sin v$\hfill .\,.\,.\,.\,.\,.\,.} & $9{,}849\,31\overset{\scriptscriptstyle\text{‘}}{0}2_n$\\[.3em]
\end{tabular}
\end{center}
\begin{center}
\begin{tabular}{@{}ll@{}}
\makebox[16em][l]{$\log\sin E\BareRoot{\frac{a}{r}}$\hfill .\,.\,.\,.\,.} & $9{,}814\,5622_n$\\[.3em]
\makebox[16em][l]{$\log\cos\frac{1}{2}\varphi$\hfill .\,.\,.\,.\,.\,.} & $9{,}996\,6567$\\[.3em]
\hline
 & $9{,}811\,2189_n$\\[.3em]
\makebox[16em][l]{$\log\sin\frac{1}{2}(v+E)$\hfill .\,.\,.} & $9{,}811\,2189_n$\\[.3em]
\end{tabular}
\end{center}
